\section{Manual (Plataforma web) - Apartado de gestión de socios}

Al ingresar a la página de socios, se visualiza en primer lugar el listado completo de todos los socios inscritos en el club. En este padrón se indica el estado actual de cada uno, el cual puede ser activo, moroso, inactivo o suspendido.

\subsection{Búsqueda, filtrado y detalle del socio}
Dentro de la pantalla principal, el sistema ofrece distintas acciones para localizar y gestionar a los usuarios:

\begin{itemize}
    \item \textbf{Búsqueda directa:} Permite buscar a un socio en particular ingresando su número (por ejemplo, el socio 1200).
    \item \textbf{Filtros:} Es posible filtrar el listado según la categoría (por ejemplo, infantiles), por su estado financiero o por las disciplinas en las que se encuentran inscritos.
\end{itemize}

Al ingresar al detalle de un socio específico, la plataforma permite acceder a su información completa:
\begin{itemize}
    \item Ver sus datos personales.
    \item Consultar las inscripciones a disciplinas y las reservas que tuviera activas.
    \item Visualizar los pagos pendientes (como las cuotas sociales) y marcarlas manualmente como pagadas en caso de que el socio realice el abono por caja.
\end{itemize}

\subsection{Creación de nuevos socios}
Desde este mismo apartado es posible dar de alta a nuevas personas mediante un formulario de creación de socios. Para completarlo, se deben ingresar los siguientes datos obligatorios y opcionales:

\begin{itemize}
    \item Nombre y apellido (por ejemplo, Axel González).
    \item Fecha de nacimiento (por ejemplo, 1 de enero del año 2000).
    \item Género.
    \item Tipo y número de documento.
    \item Correo electrónico (mail).
    \item Teléfono y dirección (datos opcionales).
\end{itemize}

Una vez confirmada la acción, el socio queda creado y el sistema le asigna automáticamente un número de socio (por ejemplo, 1304). Dado que el listado general se ordena numéricamente, el nuevo registro aparecerá posicionado al final del padrón.
